% supplementary_zh.tex — 补充材料中文版 (2026-10-16 定调: KITTI09 诚实案例 + KITTI07 新基线双报告 + 可复现性声明)
% 数字来源: se3_metrics.json / ISSUES_20261015.md 七节; 主表引用"补充材料 S.1 节"。

\section*{补充材料}
\label{sec:supplementary}

\subsection*{S.1 尺度头在 KITTI~07 上的零样本失效}

单目尺度校正头（方法章）仅在 CARLA 图像上训练一次。在真实 KITTI~07 数据上其预测发生反转：头输出 $\log s \approx -1.1$（即 $s\approx0.3$，把本已欠尺度的 VO 进一步缩小），而该序列需要约 $s\approx1.8$ 的尺度校正。将其部署——以与 CARLA 上相同的部署配置（死区 $0.30$）replay EKF、由头驱动尺度状态——会使真 VIO EKF 基线从 17.75~m 恶化到 23.13~m（Sim(3) 口径）：滤波器尺度状态从其运动学一致的值（VO 单位下约 5.5）塌向错误的深度预测（约 0.5），拖垮融合轨迹。

两个零样本参考点进一步量化了无真值时可恢复的边界：转弯运动学尺度观测器（按构造可零样本使用）在该序列收敛到 $s\approx5.65$（VO 单位），与 EKF 尺度态自行达到的值（$\approx5.54$）基本一致，而真实全局尺度为 $\approx6.78$——即残差全局尺度误差对该序列上的两个零样本尺度观测器而言均结构性不可达。因此，正文所有 KITTI 结果均在未缩放的 EKF 核之上运行闭环线，论文的度量尺度主张由四个仿真场景加（与尺度无关的）KITTI 闭环边际承载。

\subsection*{S.2 KITTI~07 基线交叉校验（双报告）}

正文的 KITTI~07 基线由我们的历史降级流水线生成（视觉里程计经同一 10 状态 EKF 与 IMU 融合，尺度误差剖面具有单目部署代表性）。由于该流水线先于最终版重新实现，我们额外在\emph{真}视觉惯性基线上评估——由当前公开流水线（\texttt{kitti\_prepare\_fusion.py}：ORB 视觉里程计经同一 EKF 与 OXTS IMU 融合）产生：

\begin{center}
\scriptsize
\setlength{\tabcolsep}{3.5pt}
\begin{tabular}{@{}lcc@{}}
\toprule
KITTI~07 配置 & Sim(3) 重拟合 & SE(3) 刚性 \\
\midrule
EKF（正文基线）          & 22.09 m & 30.45 m \\
其上的 NeuroLocMap（闭环线）     & \textbf{18.45 m} & \textbf{29.99 m} \\
\midrule
EKF（真 VIO 流水线，本节） & 17.75 m & 44.01 m \\
其上的 NeuroLocMap（闭环线）     & 23.12 m & \textbf{40.99 m} \\
\bottomrule
\end{tabular}
\end{center}

这一对比在两个方向上都有启发性。在仅形状的 Sim(3) 口径下真 VIO 基线更强（17.75~m），闭环线未能改善它（23.12~m）：Sim(3) 重拟合已吸收了该轨迹的主导误差——约 1.8 的全局尺度因子——留下闭环线的纯平移校正无法战胜的形状残差。而在刚性 SE(3) 口径下——尺度误差\emph{无法}被吸收——闭环线仍改善其基线（40.99 对 44.01~m，$-6.8\%$）。此交叉校验佐证了论文对 KITTI 结果的解读：经验地图闭环在真实数据上提供真实的全局一致性边际，且该边际的大小取决于口径允许对齐器吸收多少全局尺度误差。

\subsection*{S.3 KITTI~09：持续性门作为诚实的空操作}

为在正文五个场景之外探查门链，我们以完全相同的零样本配置处理了 KITTI~09（1591 帧，1705.1~m，大型城市环路）。结果是一次设计使然的空操作：闭环线与其 EKF 核打平（Sim(3) 42.37 对 42.34~m；SE(3) 54.63 对 54.77~m）。

该轨迹恰含一个真闭环事件——末帧回到起点区域（GT 位置差 3.8~m）。两个嵌入前端都以高置信度将该对排第一（零样本相似度 0.988，微调 0.994），但重访发生在序列最后两帧：持续性门（15 个连续候选帧，方法章）在序列结束前最多积累到 2 的连续计数，因而正确地拒绝再锚定。强制放行该事件的 what-if 仿真仅带来 $-0.12$~m（Sim(3)）——闭环来得太迟，无足轻重。与此同时，门链经朝向门拒绝了两个高相似度的虚假候选（GT 位置差 90/134~m）。地图本身健康（964 个经验节点、929 个视觉模板）。

嵌入对比还揭示一个负迁移结果：微调嵌入头——在 CARLA 城镇上提升查准率——在 KITTI~09 上把真环对排到\emph{接受门限之下}（所有阈值下查全率 0\%），而零样本嵌入将其排第一（查全率 100\%）。域偏移可以无声地退化一个在另一条真实序列（KITTI~07）上可迁移的组件；因此我们报告两种变体，并将零样本嵌入作为跨域部署的默认选择。

\subsection*{S.4 可复现性声明}

所有 CARLA 场景由公开仓库的采集器采集（智能体速度 12~km/h，20 辆自动驾驶车辆与 20 名运动学直驱行人），可由其端到端复现。KITTI~07 与 KITTI~09 由同一公开流水线处理：原始序列组装（\texttt{kitti09\_assemble\_raw.py}）后经真视觉惯性 EKF 转换为标准四文件格式（\texttt{kitti\_prepare\_fusion.py}）。KITTI~07 的两条基线并存并在 S.2 节并排报告：正文基线轨迹在发布中原样保留（\texttt{\_p11\_backup\_20261014/KITTI07/fp\_old\_synthetic.txt}），真 VIO 基线可由当前流水线从原始数据再生。论文所有指标表均由仓库中的评估脚本（\texttt{se3\_metrics.py}，双口径）从发布的轨迹文件计算，并带交叉校验以确保各表逐位一致。
